
The data columns left after the data cleaning are: \texttt{VehPower}, \texttt{VehAge}, \texttt{DrivAge}, \texttt{BonusMalus}, \texttt{VehGas}, \texttt{Density}, and these will be used as the features for our model, the variables used to teach the algorithms the necessary patterns for the prediction.\\
The columns on which we base our \texttt{Risk} calculations are the following: \texttt{ClaimNb} and \texttt{Exposure}, therefore not included in the feature set. These will steer the predictions in the correct way, this will be the information the models are to predict.

\subsection{Claims Risk}
How do we calculate it, why (idk)\\
Very imbalanced~\autoref{fig:distribution-risk}

\begin{figure} [h]
    \centering
    \caption{Histogram of Risk values}
    \label{fig:distribution-risk}
\end{figure}

%%%%%%%%%%%%%%%%%%%%%%section%%%%%%%%%%%%%%%%%%%%%%%%%%
\subsection{PCA \& Clustering}
Clustering is defined as an unsupervised learning method, which groups the observations made based on similarity, with the goal of identifying homogeneous subgroups of insurance policy claimers; telling also by the name.

For better clustering, the dimensions of the feature space should be reduced, which is why the method called PCA was introduced in this project. P(rincipal) C(omponent) A(nalysis) is a technique of dimensionality reduction which aims to turn the correlated numerical variables into a smaller set of uncorrelated ones. Where, before, the numerical features of \texttt{DrivAge}, \texttt{BonusMalus} or \texttt{Density} could be correlated, after applying the PCA, the data's high dimensionality will be stabilized and any multi-collinearity will be reduced, as the resulting set should be smaller, uncorrelated.
Further data cleaning ensued for the specific requirements of PCA, such as excluding the categorical variables and standardizing the data to prevent the dominion of large scaled values over the analysis.
Four principal components need to be used to retain 89\% of the data; which means that without significant loss most of the original information is represented in a lower dimensional space. As such, global structure is preserved and variance is recognized in fewer dimensions.

\begin{figure} [h]
    \centering
    \caption{PCA with Risk as color | Clustering on PCA data.}
    \label{fig:PCA-Risk-and-Clustering}
\end{figure}

The clustering is the following step and from various methods, we applied K-Means clustering on our data, as it's finely compatible with continuous, scaled data, resulting from the principal component analysis. The main idea behind this clustering method is the partitioning of the data into \texttt{k} clusters, \texttt{k} being the number of such groups created, through minimizing the variance between different clusters.
To find the right value for \texttt{k}, a silhouette score was used as an evaluation metric when trying different values. This score primarily measures similarity of each observation to their own cluster and the quality of the separation between clusters, after which the best one is chosen. The \texttt{k} which maximized the number of the silhouette score, is the one that is chosen for the final \texttt{k} split.
In summary of the vein of clustering, meaningful behavior and risk patterns will be displayed together, letting the visualization speak for different groups of the insurance customers.